\chapter{Brot aufbewahren}%
\label{ch:storing-bread}
\begin{quoting}
In diesem Kapitel werden verschiedene Methoden zur Aufbewahrung von Brot
besprochen, jeweils mit ihren Vor- und Nachteilen. So kann das Brot auch zu
einem spaeteren Zeitpunkt bestmoeglich genossen werden.
\end{quoting}

Eine Zusammenfassung findet sich in Tabelle~\ref{table:bread-storage}, mit
Details und Erklaerungen im weiteren Verlauf dieses Kapitels.
\begin{table}[!htb]
    \centering
        \def\arraystretch{1.6}
\begin{tabular}{@{}p{0.23\textwidth}p{0.33\textwidth}p{0.33\textwidth}@{}}
\toprule
\textbf{Methode}     & \textbf{Vorteile}            & \textbf{Nachteile}       \\
\midrule
Raumtemperatur                          & Die einfachste Methode. Am besten für Brot, das innerhalb eines Tages gegessen wird.
                                            Kruste bleibt typischerweise knusprig, wenn die Luftfeuchtigkeit nicht zu hoch ist.
                                          & Brot trocknet sehr schnell aus.\\

Raumtemperatur in luftdichtem Behälter    & Gut für bis zu eine Woche.
                                          & Brot muss getoastet werden, damit die Kruste wieder knusprig wird. Schimmelt schneller.\\

Kühlschrank                                    & Brot bleibt wochenlang gut. Kann ohne luftdichten Behälter etwas austrocknen.
                                          & Brot muss getoastet werden. Benötigt Kühlschrank und Energie.\\

Gefrierschrank                                   & Brot bleibt jahrelang gut.
                                          & Muss aufgetaut und dann getoastet werden. Benötigt Gefrierschrank und Energie.\\
\bottomrule
\end{tabular}

        \caption[Moeglichkeiten zur Brotlagerung]{Eine Tabelle mit den Vor-
            und Nachteilen verschiedener Brotlagerungsmoeglichkeiten.}%
        \label{table:bread-storage}
\end{table}

\section{Zimmertemperatur}

Die gaengigste Methode ist die Aufbewahrung des Brotes bei Zimmertemperatur.
Nach dem Abschneiden einer Scheibe das Brot mit der Schnittflaeche nach unten
lagern.

Diese Methode eignet sich gut, wenn man das Brot innerhalb eines Tages essen
moechte. Die Kruste bleibt knusprig und wird nicht weich\footnote{Je hoeher
    die Luftfeuchtigkeit im Raum, desto schneller wird die Kruste weich.}.
Der groesste Nachteil dieser Methode ist, dass das Brot schnell hart wird.
Mit der Zeit verdunstet immer mehr Wasser aus der Krume des Brotes.
Schliesslich wird das Brot sehr hart und nicht mehr essbar. Je mehr Wasser man
fuer das Brot verwendet, desto laenger bleibt es gut. Ein Rezept mit niedriger
Hydration kann nach 1--2 Tagen austrocknen; ein Brot mit hoher Hydration
braucht 3--4 Tage zum Austrocknen.

Wenn das Brot ausgetrocknet ist, kann man es etwa 10 bis 15~Sekunden unter
fliessendes Wasser halten. Dieses Wasserbad ermoeglicht es der Staerke in der
Krume, viel Wasser aufzunehmen. Anschliessend das Brot erneut im Ofen backen.
Das Ergebnis ist fast so gut wie frisch.

Eine weitere Moeglichkeit fuer ausgetrocknetes Brot ist die Herstellung von
Semmelbroeseln. Diese Brotkruemel koennen in nachfolgende Laibe eingearbeitet
werden. Sie koennen auch als Grundzutaten fuer andere Rezepte wie
\emph{Knoedel} verwendet werden\footnote{\emph{Knoedel} ist ein
    oesterreichisches Gericht, das altes Brot als Grundlage verwendet.
  Brotkruemel und altbackenes Brot werden mit Eiern vermengt, manchmal
  werden auch Spinat oder Schinken hinzugefuegt. Der Teig wird dann in
  Salzwasser gekocht.
}.

\section{Zimmertemperatur in einem Behaelter}

Wie bei der vorherigen Variante kann man das Brot auch in einem Behaelter
aufbewahren. Das kann eine Papiertuete, eine Plastiktuete oder ein
Brotkasten sein. Die Papiertuete und die meisten Brotkaesten sind nicht
vollstaendig luftdicht, sodass ein Teil der Luft aus dem Behaelter entweichen
kann. Das bedeutet auch, dass das Brot leicht austrocknet.

Bei einer versiegelten Tuete wie einer Plastiktuete behaelt das Brot viel
Feuchtigkeit. Das Brot bleibt ueber einen laengeren Zeitraum gut. Allerdings
verliert die Kruste gleichzeitig ihre Knusprigkeit. Ein Teil des Wassers
diffundiert in die Tuete und wird dann von der Kruste wieder aufgenommen. Wenn
man die knusprige Kruste moechte, ist es am besten, das Brot zu toasten.

Ein weiteres Problem bei Aufbewahrungsbehaeltern ist die natuerliche
Schimmelkontamination. In dem Moment, in dem das Brot aus dem Ofen genommen
wird, beginnt es sich mit Schimmelsporen aus der Luft zu kontaminieren. Die
Sporen sind mikroskopisch klein und ueberall vorhanden. Schimmelsporen
gedeihen am besten in einer feuchten Umgebung. Durch das Einlegen des Brotes
in einen Behaelter schafft man ein Schimmelparadies. Ein reiner Hefeteig
beginnt so innerhalb weniger Tage zu schimmeln. Sauerteigbrot bleibt laenger
gut, da die Saeure ein natuerlicher Schimmelhemmer ist.

\section{Kuehlschrank}

Nach meiner eigenen Erfahrung funktioniert die Aufbewahrung von Brot im
Kuehlschrank gut, solange man einen verschliessbaren Behaelter verwendet, auch
wenn einige Quellen sagen, dass das Brot im Kuehlschrank
austrocknet~\cite{storing+bread}. Angeblich foerdert der Kuehlschrank die
Wanderung von Fluessigkeit aus der Krume an die Brotoberflaeche.

Meiner Erfahrung nach liegt der Trick in der Verwendung eines
verschliessbaren Behaelters. Mit einem verschliessbaren Zip-Beutel bleibt die
ueberschuessige Feuchtigkeit im Beutel und stellt sicher, dass das Brot nicht
so schnell austrocknet. Bei Zimmertemperatur wuerde das zum Schimmeln fuehren.
Bei niedrigeren Temperaturen kann das Brot so wochenlang gut bleiben. Die
Kruste verliert jedoch ihre Knusprigkeit, weshalb Toasten empfohlen wird.

\section{Einfrieren}

Eine weitere grossartige Moeglichkeit fuer die Langzeitlagerung ist das
Einfrieren. Den ganzen Laib in Scheiben schneiden und Portionen erstellen, die
man innerhalb eines Tages verbrauchen kann. Jede Portion in einen separaten
Behaelter legen und ins Gefrierfach stellen.

Wenn man frisches Brot essen moechte, morgens einen der Behaelter oeffnen und
das Brot einige Stunden auftauen lassen. Das ist noetig, damit man die
zusammengefrorenen Scheiben leicht trennen kann. Die Scheiben im Toaster oder
im Ofen toasten, bis sie die gewuenschte Knusprigkeit haben.

Diese Variante ist hervorragend fuer die Langzeitlagerung. Persoenlich habe
ich gerne ein paar Scheiben Brot als Notfallreserve eingefroren, wenn ich
keine Zeit zum Backen hatte.

Eine Studie aus dem Jahr 2008 deutet darauf hin, dass das Einfrieren und
Toasten von Brot gesundheitliche Vorteile haben koennte. Dadurch koennten die
Staerkemolekuele resistenter gegen die Verdauung werden und so die
Blutzuckerreaktion des Koerpers um fast \qty{40}{\percent}
senken~\cite{freezing+toasting+bread}.
